\section{Introduction}

The incident edges of a vertex determine a natural likelihood classifier
when the community parameters and the other labels are known. A spectral
method replaces those unknown labels by a low-dimensional representation
of the graph. We ask whether the degrees and a single normalized
eigensystem retain enough information to match the error of the
incident-row likelihood classifier up to a constant factor.

We construct an explicit classifier with this property for stochastic
block models with arbitrary fixed \(K\). Its normalization uses an
empirical positive degree floor. The decoder retains the raw degrees,
reconstructs a signed matrix with exactly those row sums, and alternates
symmetric Bernoulli block fitting with ordinary likelihood classification.
The only initialization is a deterministic radius-count procedure on
normalized spectral rows. At least one likelihood update is required;
any finite number of the specified updates is allowed.

The main result holds for every diverging degree scale \(d=np\), subject
to the upper sparsity condition \(np^2\le C_1\). In particular, it covers
\(d=\log\log n\) and slower divergence. The proof controls average
classification risk directly, allowing isolated vertices and exceptional
spectral rows. A graph-wide bound on the largest spectral row would
discard this range. Instead, a truncated-energy estimate bounds the
probability that the tested incident row actually encounters an
exceptional coefficient.

The target comparison is sharp at the level of risk. For a pair of
communities with exact Bernoulli Chernoff information \(I\), the oracle
error is of order \(d^{-1/2}e^{-I}\). Sharp Chernoff bounds in community
detection are developed by Zhou and Li~\cite{zhouli}; spectral error
exponents are studied by Zhang~\cite{zhang}. Here the local spectral
and fitted-likelihood errors are bounded on the scale needed to
preserve a constant risk multiplier. The geometric input uses sparse
graph regularization~\cite{llv2017} and, at larger degree scales,
a bounded-entry matrix norm inequality~\cite{bvh}.

\section{Model and the incident-row oracle}\label{sec:statements}

There are \(n\) vertices and a fixed number \(K\ge2\) of communities.
Let \(z\in[K]^n\), let \(Z\) be its indicator matrix, and let
\(n_a=\#\{i:z_i=a\}\). The symmetric adjacency matrix has zero diagonal
and independent upper-triangular entries satisfying
\begin{equation}\label{eq:mean-model}
 A_{ij}\sim\Ber(P_{z_i z_j}),\qquad
 P=pC,\qquad d=np,\qquad \Omega=ZPZ^\top,\qquad i<j.
\end{equation}
The structural constants obey
\begin{equation}\label{eq:structural}
 n_a\ge\pi_-n,\qquad C=C^\top,\qquad
 0<c_-\le C_{ab}\le c_+<\infty,\qquad
 \sigma_{\min}(C)\ge\kappa>0.
\end{equation}
The number \(K\), \(\pi_-\), \(c_-\), \(c_+\) and \(\kappa\) are fixed.
The matrix \(C\) may vary with \(n\) inside this class. Constants in
bounds may depend on these fixed quantities and on the stated upper
bound for \(np^2\), but not on \(n\). The completed mean \(\Omega\)
includes a diagonal solely for algebra; the actual mean is
\(\Omega-\diag(\Omega_{ii})\). We keep this \(O(p)\) difference in the
proof.

For a partition \(\widehat z\), define
\[
 \ell(\widehat z,z)=\frac1n\min_{\sigma\in\mathfrak S_K}
       \sum_i\ind{\sigma(\widehat z_i)\ne z_i}.
\]
The oracle knows \(P\) and all labels except the tested label. At vertex
\(i\) it maximizes the exact Bernoulli log likelihood
\begin{equation}\label{eq:sbm-oracle}
 L_{i,*}(a)=\sum_{j\ne i}
 \{A_{ij}\log P_{az_j}+(1-A_{ij})\log(1-P_{az_j})\}.
\end{equation}
It uses uniform candidate weights and a fixed rule for ties. An exact
global community-size constraint is not used to infer the missing
label. If \(z_{i,*}\) denotes this local classifier, its average
directional error is
\begin{equation}\label{eq:oracle-risk}
 R_{\mathrm{SBM},*}(n)=\frac1n\sum_i\PP\{z_{i,*}\ne z_i\}.
\end{equation}
This is the oracle used throughout the article.

\section{A decoder from degrees and normalized spectrum}
\label{restore:algorithm}

The input is \(A\) and \(K\). Fix sufficiently small constants
\(a_0,c_0>0\), depending only on the structural class. All choices and
ties below are deterministic. A fixed fallback partition defines the
algorithm whenever a required operation is undefined.

\paragraph{Encoding.}
Write
\[
 \mathbf d=A\one,\qquad D_i=\mathbf d_i,\qquad
 \bar D=\frac1n\sum_iD_i.
\]
\[
 G=\diag\{\max(D_i,a_0\bar D)\},\qquad
 S_G=G^{-1/2}AG^{-1/2}.
\]
The all-zero graph has the fixed output. Otherwise \(G\) is positive.
Retain the raw degrees and the \(K\) eigenpairs of \(S_G\) with largest
absolute eigenvalues:
\[
 U^\top U=I_K,\qquad S_GU=U\Lambda.
\]
Only \((\mathbf d,U,\Lambda)\) are used after this step. Flooring does
not require a degree vector to be a retained eigenvector.

\paragraph{Exact-degree reconstruction.}
Put
\[
 H=G^{-1/2}UU^\top G^{1/2},\qquad
 R_0=G^{1/2}U\Lambda U^\top G^{1/2}=AH.
\]
Let \(v=\mathbf d-R_0\one\), let \(R_1=R_0+v\one^\top/n\),
and let \(r\) be the vector of diagonal entries of \(R_1\). Define
\begin{equation}\label{restore:reconstruction}
 R^h=R_1+\frac{r\one^\top}{n-1}
              -\frac{n}{n-1}\diag(r).
\end{equation}
Then
\[
 R^h\one=\mathbf d,\qquad \diag(R^h)=0.
\]
The matrix \(R^h\) can be signed and need not be symmetric. Its entries
are not clipped. The formula is computed entirely from the retained
quantities. In particular, an isolated vertex has a zero reconstructed
row and does not stop classification of other vertices.

\paragraph{Initialization.}
Normalize each nonzero row of \(U\) to unit length; assign a fixed unit
direction to a zero row. Among the remaining rows, select one having
the largest number of remaining rows within distance \(1/8\), and
remove all rows within distance \(1/4\) of it. Repeat for \(K\) centers,
then assign every row to its nearest center. If fewer than \(K\)
centers can be selected, use the fixed fallback. This is a finite
procedure using \(O(Kn^2)\) distance comparisons.

\paragraph{Block fitting and classification.}
Given current labels \(\widehat z\), for \(a\le b\) let
\(\mathcal I_{ab}(\widehat z)\) contain all ordered pairs \(i\ne j\)
whose unordered label pair is \(\{a,b\}\). Let
\[
 E_{ab}=\sum_{(i,j)\in\mathcal I_{ab}(\widehat z)}R^h_{ij},
 \qquad N_{ab}=|\mathcal I_{ab}(\widehat z)|,
 \qquad \varepsilon=c_0\bar D/n.
\]
The symmetric fitted matrix is
\begin{equation}\label{eq:block-fit}
 \widehat P_{ab}=\widehat P_{ba}
   =\operatorname{clip}_{[\varepsilon,\,1/2]}(E_{ab}/N_{ab}).
\end{equation}
This maximizes the ordered-pair Bernoulli working criterion with a
symmetric block probability. The fit is interior on the event used in
the proof. Update all labels simultaneously by maximizing
\begin{equation}\label{eq:fitted-row-score}
 \widehat L_i(a)=\sum_{j\ne i}
 \{R^h_{ij}\log\widehat P_{a\widehat z_j}
   +(1-R^h_{ij})\log(1-\widehat P_{a\widehat z_j})\}.
\end{equation}
In particular, the full nonedge term is retained. There is no target
degree parameter to profile in this ordinary SBM algorithm.
Write \(\widehat L_i^{ab}=\widehat L_i(a)-\widehat L_i(b)\)
for a fitted pairwise contrast.

Perform at least one block fit and likelihood update. The output may
follow any finite number of these hard refits, including finite
data-dependent stopping. Empty fitted groups or undefined operations
use the stated fallback; the proof shows they do not occur on the
aggregate event. No leave-one-out graph, true label or true model
parameter is an algorithm input.

\section{Oracle-risk guarantee}\label{sec:main-result}

\begin{theorem}\label{restore:sbm-theorem}
Under \eqref{eq:mean-model}--\eqref{eq:structural}, suppose
\[
 d=np\longrightarrow\infty,\qquad np^2\le C_1<\infty.
\]
For sufficiently small fixed \(a_0,c_0>0\), the algorithm in
\cref{restore:algorithm}, after any finite positive number of the
specified likelihood updates, satisfies
\begin{equation}\label{eq:target}
 \E\ell(\widehat z,z)\le C R_{\mathrm{SBM},*}(n)
\end{equation}
for all sufficiently large \(n\). The same bound holds for finite
data-dependent stopping. The constant \(C\) is independent of \(n\)
and of the number of refits.
\end{theorem}

There is no condition of the form \(np\ge c\log n\).
The assumptions allow general symmetric positive block profiles and
comparable, unequal community sizes. They do not require a planted
partition or equal expected degrees. The condition \(np^2\le C_1\)
is retained as an upper sparsity assumption.

For the exact pairwise Bernoulli information \(I_{i,ab}\) defined in
\eqref{eq:sbm-information}, \cref{lem:sharp} gives
\[
 R_{\mathrm{SBM},*}(n)\asymp
 \frac{1}{n\sqrt d}\sum_i\sum_{b\ne z_i}e^{-I_{i,z_i b}}.
\]
Thus \eqref{eq:target} preserves the sharp prefactor as well as the
exponent. It asserts a bounded risk ratio; it does not assert that
the ratio tends to one. Since a nonzero loss is at least \(1/n\),
the theorem also implies exact recovery whenever
\(nR_{\mathrm{SBM},*}(n)\to0\).

\subsection{Proof structure}

The proof has four ingredients. First, degree-floor geometry gives a
fixed informative spectral gap, \(O(d^{-1/2})\) subspace error and
\(O(n)\) weighted degree energy. These aggregate bounds suffice for
the deterministic initialization and an invariant \(O(n/d)\)
misclassification basin.

Second, a same-threshold truncated-energy recursion controls large
spectral rows in expectation. By taking a fixed row threshold large
enough, the average probability of an exceptional tested row, or of
an incident edge hitting an exceptional training row, is at most
\(C_He^{-Hd}\) for any prescribed fixed \(H\). This estimate replaces
a simultaneous row bound.

Third, on the local regular event, the exact block fit has error
\(O(1/n)\), and every fitted row contrast differs by \(O(1)\) from its
true block contrast. A scalar leave-one-out identity replaces the
full-graph spectral action by a function of independent training
data and the tested incident row, again with \(O(1)\) error. Under
the oracle Chernoff tilt, the resulting coefficient error has bounded
mean and quadratic mass.

Finally, a sharp Bernoulli transfer preserves
\(d^{-1/2}e^{-I_{i,ab}}\). Exceptional coefficients are paid for only
when the incident row hits them. Polynomially accurate regularized
geometry suffices below an internal logarithmic proof threshold;
an exponentially accurate norm event supplies the complementary
range. These estimates together cover every \(d\to\infty\).
The appendices give the complete argument and the proof of
\cref{restore:sbm-theorem}.

\paragraph{Organization of related material.}
The companion notes contain degree correction, alternative earlier
decoders, counterexamples and supplementary derivations. None of
those notes is required for this article's theorem or proof.

